\documentclass[tikz,border=3mm,multi=tikzpicture]{standalone}
\usepackage{eleve-matematica-ef1}

\newcommand{\Fulcro}[2]{%
  \fill[matemazulclaro] ({#1-0.62},{#2-1.05}) --
    ({#1+0.62},{#2-1.05}) -- (#1,#2) -- cycle;
  \draw[matematica contorno] ({#1-0.62},{#2-1.05}) --
    ({#1+0.62},{#2-1.05}) -- (#1,#2) -- cycle;
}

\begin{document}

% FIGURA 1 — fig-01-balanca-equilibrada-e-desequilibrada
\begin{tikzpicture}[matematica figura, x=1cm, y=1cm]
  \path[use as bounding box] (-0.10,-0.10) rectangle (10.60,11.35);

  % Equilíbrio
  \draw[matematica contorno] (1.55,8.85) -- (9.05,8.85);
  \draw[matematica contorno] (2.25,8.85) -- (2.25,7.55);
  \draw[matematica contorno] (8.35,8.85) -- (8.35,7.55);
  \draw[matematica prato] (0.95,7.55) -- (3.55,7.55);
  \draw[matematica prato] (7.05,7.55) -- (9.65,7.55);
  \Fulcro{5.30}{8.85}
  \foreach \x/\y in {1.75/7.93,2.25/8.20,2.75/7.93}
    \fill[matemalaranja] (\x,\y) circle (0.23);
  \fill[matemaverde] (8.35,8.02) ellipse (0.48 and 0.34);
  \node[matematica rotulo] at (2.25,6.92) {3 laranjas};
  \node[matematica rotulo] at (8.35,6.92) {1 mamão};
  \node[font=\Large\bfseries,text=matemaverde] at (5.30,10.15)
    {massas iguais};

  % Desequilíbrio
  \draw[matematica contorno] (1.55,3.25) -- (9.05,4.65);
  \draw[matematica contorno] (2.25,3.38) -- (2.25,1.95);
  \draw[matematica contorno] (8.35,4.52) -- (8.35,3.10);
  \draw[matematica prato] (0.95,1.95) -- (3.55,1.95);
  \draw[matematica prato] (7.05,3.10) -- (9.65,3.10);
  \Fulcro{5.30}{3.95}
  \foreach \x/\y in {1.58/2.31,2.02/2.56,2.46/2.31,2.90/2.56,1.80/2.82,2.68/2.82}
    \fill[matemalaranja] (\x,\y) circle (0.21);
  \fill[matemaverde] (8.35,3.55) ellipse (0.48 and 0.34);
  \node[matematica rotulo] at (2.25,1.30) {mais peso};
  \node[matematica rotulo] at (8.35,2.48) {mesma massa};
  \node[font=\Large\bfseries,text=matemavermelho] at (5.30,5.50)
    {um lado mudou};
\end{tikzpicture}

% FIGURA 2 — fig-02-dois-membros-da-igualdade
\begin{tikzpicture}[matematica figura, x=1cm, y=1cm]
  \path[use as bounding box] (-0.10,-0.10) rectangle (10.60,6.95);

  \draw[matematica contorno] (1.55,4.25) -- (9.05,4.25);
  \draw[matematica contorno] (2.25,4.25) -- (0.95,2.75);
  \draw[matematica contorno] (2.25,4.25) -- (3.55,2.75);
  \draw[matematica contorno] (8.35,4.25) -- (7.05,2.75);
  \draw[matematica contorno] (8.35,4.25) -- (9.65,2.75);
  \draw[matematica prato] (0.95,2.75) -- (3.55,2.75);
  \draw[matematica prato] (7.05,2.75) -- (9.65,2.75);
  \Fulcro{5.30}{4.25}

  \node[matematica valor] at (2.25,3.30) {$3+4$};
  \node[matematica valor] at (8.35,3.30) {$7$};
  \node[font=\Huge\bfseries,text=matemalaranja] at (5.30,5.25) {$=$};
  \node[font=\large\bfseries,text=matemacinza] at (2.25,1.90)
    {primeiro membro};
  \node[font=\large\bfseries,text=matemacinza] at (8.35,1.90)
    {segundo membro};
\end{tikzpicture}

% FIGURA 3 — fig-03-mesma-adicao-nos-dois-membros
\begin{tikzpicture}[matematica figura, x=1cm, y=1cm]
  \path[use as bounding box] (-0.10,-0.10) rectangle (10.60,11.10);

  % Antes
  \node[font=\Large\bfseries,text=matemacinza] at (5.30,10.42) {antes};
  \draw[matematica contorno] (1.55,8.55) -- (9.05,8.55);
  \draw[matematica contorno] (2.25,8.55) -- (0.95,7.25);
  \draw[matematica contorno] (2.25,8.55) -- (3.55,7.25);
  \draw[matematica contorno] (8.35,8.55) -- (7.05,7.25);
  \draw[matematica contorno] (8.35,8.55) -- (9.65,7.25);
  \draw[matematica prato] (0.95,7.25) -- (3.55,7.25);
  \draw[matematica prato] (7.05,7.25) -- (9.65,7.25);
  \Fulcro{5.30}{8.55}
  \node[matematica valor] at (2.25,7.82) {$9$};
  \node[matematica valor] at (8.35,7.82) {$9$};

  % Depois
  \node[font=\Large\bfseries,text=matemalaranja] at (5.30,5.05)
    {depois de somar 3 aos dois lados};
  \draw[matematica contorno] (1.55,3.20) -- (9.05,3.20);
  \draw[matematica contorno] (2.25,3.20) -- (0.95,1.90);
  \draw[matematica contorno] (2.25,3.20) -- (3.55,1.90);
  \draw[matematica contorno] (8.35,3.20) -- (7.05,1.90);
  \draw[matematica contorno] (8.35,3.20) -- (9.65,1.90);
  \draw[matematica prato] (0.95,1.90) -- (3.55,1.90);
  \draw[matematica prato] (7.05,1.90) -- (9.65,1.90);
  \Fulcro{5.30}{3.20}
  \node[matematica valor] at (2.25,2.47) {$9+3$};
  \node[matematica valor] at (8.35,2.47) {$9+3$};
  \node[font=\Large\bfseries,text=matemaverde] at (5.30,0.52)
    {$12=12$};
\end{tikzpicture}

% FIGURA 4 — fig-04-igualdade-verdadeira-e-falsa
\begin{tikzpicture}[matematica figura, x=1cm, y=1cm]
  \path[use as bounding box] (-0.10,-0.10) rectangle (10.60,11.35);

  % Verdadeira
  \node[font=\Large\bfseries,text=matemaverde] at (5.30,10.35)
    {igualdade verdadeira};
  \draw[matematica contorno] (1.55,8.55) -- (9.05,8.55);
  \draw[matematica contorno] (2.25,8.55) -- (0.95,7.25);
  \draw[matematica contorno] (2.25,8.55) -- (3.55,7.25);
  \draw[matematica contorno] (8.35,8.55) -- (7.05,7.25);
  \draw[matematica contorno] (8.35,8.55) -- (9.65,7.25);
  \draw[matematica prato] (0.95,7.25) -- (3.55,7.25);
  \draw[matematica prato] (7.05,7.25) -- (9.65,7.25);
  \Fulcro{5.30}{8.55}
  \node[matematica valor] at (2.25,7.82) {$12$};
  \node[matematica valor] at (8.35,7.82) {$12$};

  % Falsa
  \node[font=\Large\bfseries,text=matemavermelho] at (5.30,5.55)
    {igualdade falsa};
  \draw[matematica contorno] (1.55,3.25) -- (9.05,4.65);
  \draw[matematica contorno] (2.25,3.38) -- (0.95,1.95);
  \draw[matematica contorno] (2.25,3.38) -- (3.55,1.95);
  \draw[matematica contorno] (8.35,4.52) -- (7.05,3.10);
  \draw[matematica contorno] (8.35,4.52) -- (9.65,3.10);
  \draw[matematica prato] (0.95,1.95) -- (3.55,1.95);
  \draw[matematica prato] (7.05,3.10) -- (9.65,3.10);
  \Fulcro{5.30}{3.95}
  \node[matematica valor] at (2.25,2.50) {$9$};
  \node[matematica valor] at (8.35,3.65) {$8$};
  \node[font=\Large\bfseries,text=matemavermelho] at (5.30,0.45)
    {$9\neq8$};
\end{tikzpicture}

\end{document}
